% Page budget: 1.55 pages | Owns: augmentation topology, additional states, encoder, and closed-loop objective.
% WRITING GUIDE
% Purpose: Define the final augmentation architecture and explain how it is trained inside the known feedback loop.
% Start here: Current implementation decisions D-129, D-140, D-141, D-180, D-220, D-233, and D-237 in docs/decisions.md.
% Supporting material: Quinten's 18 September outline feedback, the Meeting-audit synthesis, Thesis-writeup/Documentation/TRAINING-DESIGN.md, and docs/writeup figures.
% Mandatory papers: Read Hoekstra's 2025 LFR augmentation paper, the 2026 first-principles augmentation paper, the encoder-initialization paper, and Kessels' Chapter 5 before writing this section.
% Research-plan anchor: Preserve Aspect 2's dynamic parallel augmentation objective; document the departure from open-loop training and earlier state routing.
% Current decisions: Separate physical and additional states, keep the controller outside the model, score the complete training window without burn-in, and use one network for the physical correction and added-state update.
% Choices to justify: Final reported routing, added-state dimension, ANN width and depth, zero initialization, encoder history, closed-loop objective, rollout length, full-window scoring, and validation horizon. Use the configuration of the reported thesis run, not a superseded routing experiment.
% Horizon check: Treat encoder history, scored training window, joint-estimation horizon, and free-run validation horizon separately; relate candidates to relevant stable modes and test sensitivity to slower dynamics and free-integrator accumulation.
% Implementation audit: Read scripts/gantry/gantry_interconnect_dynamic.py, gantry_dynamic/{model,config,controller,data,training}.py, and model_augmentation/fit_systems/{interconnect,closed_loop,pre_encoder}.py.
% Reference check: Verify sources for parallel LFR augmentation, encoder initialization, closed-loop state extension, and any claimed architectural precedent against the original paper or thesis.
% Cautions: Earlier routing plans are superseded; distinguish the truth-only hidden absorber from learned states; closed-loop fit alone does not prove autonomous plant recovery.
% Figures: Absorber schematic, author's updated architecture, and thesis-owned common-reference loops. No residual panel or horizon figure.
% Section is finished when: Every signal has a defined frame and timing, the RK4 boundary is explicit, and the described architecture matches the final code.
%
% SOURCE MAP
% Hoekstra et al. 2025 EJC: dynamic-parallel topology, additional states, encoder-based truncated simulation.
% Hoekstra et al. 2026 Automatica submission: extended LFR formulation, baseline-preserving model/encoder initialization, joint identification context.
% Hoekstra et al. 2026 encoder paper: reconstructability-based W^b initialization; it does not identify W^a or test dynamic augmentation.
% Kessels 2025 thesis: direct closed-loop rollout, controller equations, and reconstruction of the machine controller state for each window.
% Drenth 2025 thesis: LPV-specific augmentation precedent; not the unrestricted MLP structure used here.
% This project: gantry specialization, residual-loop implementation, exact routing, RK4 boundary, and verification.
\section{Dynamic LPV-LFR Augmentation}\label{sec:augmentation}

The baseline of Section~\ref{sec:lfr} does not describe every effect of the
machine.
Garc\'ia-Herreros et al.\ neglect the resonance of the supporting base, the
vibration of the cross-arm on its flexible plates, the detent forces of the
linear motors and the variation of friction along the
stroke~\cite[Sec.~2.4]{garcia2013model}, and \eqref{eq:eom} also omits
their Coulomb friction.
We augment the baseline with a learned dynamic component and estimate the
physical parameters $\theta_{\mathrm{base}}$, the network parameters
$\theta_{\mathrm{aug}}$ and the encoder parameters $\eta$ jointly from
records measured under feedback.
Section~\ref{sec:aug_structure} defines the augmented model,
Section~\ref{sec:aug_encoder} the initial state of each training window, and
Section~\ref{sec:aug_closed_loop} the closed-loop objective.

\subsection{Augmented model}\label{sec:aug_structure}
An omitted vibration mode adds model order.
Refitting $\theta_{\mathrm{base}}$ cannot supply it, and neither can a
static correction of the state update, because both act only on the
existing states; dynamic augmentations add states for this
purpose~\cite[Sec.~2]{hoekstra2025lfr}.
We therefore identify the dynamic-parallel
structure~\cite[Table~1]{hoekstra2026lfr}
\begin{equation}
\begin{aligned}
 \tilde x_{k+1}
 &=f_{\mathrm{base}}(\theta_{\mathrm{base}},\tilde x_k,u_k)
   +f_{\mathrm{aug}}(\theta_{\mathrm{aug}},\hat x_k,u_k),\\
 \bar x_{k+1}
 &=g_{\mathrm{aug}}(\theta_{\mathrm{aug}},\hat x_k,u_k),\\
 \hat y_k&=h_{\mathrm{base}}(\tilde x_k)
   =\begin{bmatrix}P^\top & 0\end{bmatrix}\tilde x_k ,
\end{aligned}
\label{eq:aug_transition}
\end{equation}
where $\tilde x=\col(q,\dot q)\in\R^6$ is the physical state of
Section~\ref{sec:lfr}, $\bar x\in\R^{n_{x_a}}$ collects the additional
states, $\hat x=\col(\tilde x,\bar x)$, $u_k=Pu_{\mathrm{act},k}$ is the
generalized force of \eqref{eq:Ptransform}, and $\hat y_k$ predicts the
measured positions $y_k=q_{\mathrm{act},k}$.
As in the dynamic-parallel structure, $f_{\mathrm{aug}}$ acts on every
physical-state row, the positions included.
The output map is not learned, so the additional states reach $\hat y$
only through $f_{\mathrm{aug}}$; they have no assigned physical meaning and
are not the states of the absorber in the simulated truth of
Section~\ref{sec:setup}.

$f_{\mathrm{base}}$ is one RK4 step of \eqref{eq:eom} over the sample
period $T_s$, with the input held over the step, as
in~\cite[Sec.~7.2]{hoekstra2026lfr}.
Each stage evaluates $M(Y)$ at its own state, so $f_{\mathrm{base}}$
remains self-scheduled.
All signals are normalised with the statistics of the training
records~\cite[Sec.~5.3]{hoekstra2026lfr}.
The learned functions act on the result of the RK4 step, not inside it:
$f_{\mathrm{base}}$ integrates the continuous-time physics, and
$f_{\mathrm{aug}}$ is a discrete-time correction.
The LFR of Section~\ref{sec:lfr} is therefore unchanged, whereas the
LPV-LFR augmentation of Drenth adds states and latent variables to the
baseline LFR itself~\cite[Sec.~5.2]{drenth2025thesis}.

One multilayer perceptron with tanh hidden layers supplies both learned
functions,
\begin{equation}
 \begin{bmatrix}
  f_{\mathrm{aug}}(\theta_{\mathrm{aug}},\hat x_k,u_k)\\
  g_{\mathrm{aug}}(\theta_{\mathrm{aug}},\hat x_k,u_k)
 \end{bmatrix}
 =\phi_{\mathrm{aug}}\bigl(\theta_{\mathrm{aug}},\col(\hat x_k,u_k)\bigr),
 \label{eq:aug_mlp}
\end{equation}
where $\theta_{\mathrm{aug}}$ collects its weights and biases, the first
six outputs form $f_{\mathrm{aug}}$, and the last $n_{x_a}$ form
$g_{\mathrm{aug}}$ (Fig.~\ref{fig:aug_structure}).

\begin{figure}[t]
 \centering
 \resizebox{\columnwidth}{!}{\input{figures/tex/augmentation_structure}}
 \caption{Augmented model \eqref{eq:aug_transition}. The physics step
 $f_{\mathrm{base}}$ and the network $\phi_{\mathrm{aug}}$ act in parallel
 on $\hat x_k$; $\phi_{\mathrm{aug}}$ supplies $f_{\mathrm{aug}}$ and
 $g_{\mathrm{aug}}$, and the output map is not learned
 ($h_{\mathrm{aug}}=0$). The encoder $\psi$ sets $\hat x$ at each window
 start. During training, $u_k$ is the closed-loop input $\hat u_k$ of
 \eqref{eq:aug_controller}.}
 \label{fig:aug_structure}
\end{figure}

The augmented model is well posed.
Within one step no learned function feeds another, so the computational
graph is acyclic; with $f_{\mathrm{base}}$ and the tanh network $C^2$,
well-posedness follows from~\cite[Thm.~9]{hoekstra2026lfr}.
$f_{\mathrm{base}}$ is smooth because $M(Y)\succ0$ for every $Y$ at every
iterate (Sections~\ref{sec:lfr} and~\ref{sec:params}), so no value of $Y$
written by $f_{\mathrm{aug}}$ makes the baseline ill posed.

Training starts from the baseline.
The output layer of the network, with weights $W_L$ and bias $b_L$, is
initialised at zero,
\begin{equation}
 W_L=0,\quad b_L=0
 \quad\Longrightarrow\quad
 f_{\mathrm{aug}}=0,\quad g_{\mathrm{aug}}=0 ,
 \label{eq:aug_zero_init}
\end{equation}
so the augmented model reproduces the baseline from the same physical
initial state, the initialisation condition
of~\cite[Eq.~(29)]{hoekstra2026lfr}.

From this start, $f_{\mathrm{aug}}$ can also learn changes that a different
$\theta_{\mathrm{base}}$ would produce, so the split between the two
components is not unique~\cite[Sec.~5.2]{hoekstra2026lfr}.
We do not use the parameter regularisation
of~\cite[Eq.~(26)]{hoekstra2026lfr}: it bounds the drift of
$\theta_{\mathrm{base}}$ but does not remove this ambiguity.
Section~\ref{sec:obc} constrains the split instead.

\subsection{Encoder and initial state}\label{sec:aug_encoder}
Each training window simulates \eqref{eq:aug_transition} from an initial
state, but only the positions are measured.
Following Beintema et al.~\cite{beintema2023subnet}, an encoder estimates
this state from the input and output history,
\begin{equation}
\begin{aligned}
 \hat x_{\tau|\tau}&=\psi_\eta(\xi_\tau)
 =\begin{bmatrix}W^b\\ W^a\end{bmatrix}\xi_\tau+\tilde\psi(\xi_\tau),\\
 \xi_\tau&=\col\bigl(y_{\tau-n:\tau},\,u_{\mathrm{act},\tau-n:\tau}\bigr),
\end{aligned}
\label{eq:aug_encoder}
\end{equation}
where $\tau$ is the window start, $n$ the encoder history, $W^b$ and $W^a$
the linear maps to the physical and additional states, and $\tilde\psi$ an
MLP, as in~\cite[Eq.~(8)]{hoekstra2026encoder}.
Unlike~\cite[Eq.~(8)]{hoekstra2026encoder}, the history includes the
current sample, which the reconstructability map
uses~\cite[Eq.~(15)]{hoekstra2026encoder}.

$W^b$ is initialised from the baseline.
Linearising \eqref{eq:eom} about rest at $Y=0$ with the nominal parameters
of Appendix~\ref{app:params} gives
\begin{equation}
 A_c=\begin{bmatrix}0 & I\\ -M_0^{-1}K & -M_0^{-1}C\end{bmatrix},\qquad
 B_c=\begin{bmatrix}0\\ M_0^{-1}P\end{bmatrix},\qquad
 C_d=\begin{bmatrix}P^\top & 0\end{bmatrix},
\label{eq:aug_lin}
\end{equation}
with $M_0=M(0)$ from Section~\ref{sec:lfr}.
Its zero-order-hold discretisation $(A_d,B_d,C_d)$ at $T_s$, normalised as
the signals, gives $W^b$ as the reconstructability
map~\cite[Eqs.~(16), (17)]{hoekstra2026encoder}
\begin{equation}
 W^b=\begin{bmatrix}
  A_d^nO_n^{+} & \;r_n-A_d^nO_n^{+}T_n
 \end{bmatrix},
 \label{eq:aug_wb}
\end{equation}
where $O_n$ is the observability matrix of $(A_d,C_d)$ over $n+1$ samples,
$T_n$ the Toeplitz matrix of its input-to-output responses, $r_n$ the
input-to-state map over the same samples, and the two column blocks act on
the output and input parts of $\xi_\tau$.
All positions are measured, so $(A_d,C_d)$ is observable and $O_n^{+}$ is a
left inverse~\cite[Sec.~3.1]{hoekstra2026encoder}.
This sets $W^b$ analytically, where~\cite[Eq.~(30)]{hoekstra2026lfr} fits
the baseline encoder to simulated baseline states.

The rows $W^a$ have no baseline counterpart and are drawn randomly, and the
output layer of $\tilde\psi$ starts at zero.
By \eqref{eq:aug_zero_init}, the encoded additional states do not affect
the output at initialisation.
The encoder is trained jointly with the model, so the linearisation only
sets the starting value of $W^b$.

\subsection{Closed-loop rollout and objective}\label{sec:aug_closed_loop}
The records are measured under feedback, and the model is used to predict
the machine under feedback: its response to a given reference, feedforward
and controller.
We therefore fit this closed-loop response, rather than the open-loop
simulation driven by the recorded input of~\cite[Eq.~(22)]{hoekstra2026lfr}.

The baseline gives a second reason.
$X$ and $Y$ have no stiffness in \eqref{eq:damping_stiffness_matrices}, so
in an open-loop rollout a velocity error integrates into a position error
that persists.
Inside the loop, the controller acts on such errors as it does on the
machine, provided the model in feedback with the controller is stable.
At initialisation the learned functions are zero, so this holds whenever
the controller stabilises the baseline; Kessels starts from near-zero
network weights for the same reason~\cite[Sec.~5.2.3]{kessels2025ai}.

Kessels trains an augmented model with the known feedback controller inside
each truncated window~\cite[Eqs.~(5.12), (5.13c), (5.13d)]{kessels2025ai}
and keeps the controller outside the learned model, so that it can be
replaced in evaluation~\cite[Sec.~5.3.2.3]{kessels2025ai}.
We adopt both.
Unlike~\cite[Eq.~(5.12)]{kessels2025ai}, which optimises the network
weights only, we estimate $\theta_{\mathrm{base}}$ jointly.

The recorded loop is driven by a reference $r$ and a feedforward
$u^{\mathrm{ff}}$ through a linear controller $(A_K,B_K,C_K,D_K)$ with
state $s$, which maps the position error in $\qact$ to actuator forces.
The model loop receives the same $r$, $u^{\mathrm{ff}}$ and controller, and
differs only in its output $\hat y$ (Fig.~\ref{fig:aug_closed_loop}):
\begin{equation}
\begin{aligned}
 s_{k+1}&=A_K s_k+B_K(r_k-y_k), \\
 u_{\mathrm{act},k}&=C_K s_k+D_K(r_k-y_k)+u^{\mathrm{ff}}_k,\\
 \hat s_{k+1}&=A_K \hat s_k+B_K(r_k-\hat y_k), \\
 \hat u_{\mathrm{act},k}&=C_K \hat s_k+D_K(r_k-\hat y_k)+u^{\mathrm{ff}}_k .
\end{aligned}
\label{eq:aug_direct_controller}
\end{equation}
Subtracting the recorded loop from the model loop removes $r$ and
$u^{\mathrm{ff}}$,
\begin{equation}
\begin{aligned}
 x^c_{k+1}&=A_K x^c_k+B_K(y_k-\hat y_k),\qquad x^c_\tau=0,\\
 \hat u_k&=P\bigl(u_{\mathrm{act},k}+C_K x^c_k+D_K(y_k-\hat y_k)\bigr),
\end{aligned}
\label{eq:aug_controller}
\end{equation}
where $x^c=\hat s-s$ and $\hat u_k=P\hat u_{\mathrm{act},k}$ is the input
of \eqref{eq:aug_transition}.
This form is exact when both loops use the same linear controller, not
scheduled on the model state, with the same feedforward and without
saturation.
When the machine's controller runs at another rate than $T_s$, its
discretisation at $T_s$ is the one approximation.

\begin{figure}[t]
 \centering
 \includegraphics[width=\columnwidth]{closed_loop_same_reference.pdf}
 \caption{Recorded plant loop (top) and augmented-model loop (bottom),
 driven by one shared reference $r_k$ and feedforward $u_k^{\mathrm{ff}}$,
 with the same feedback controller. The model output is evaluated
 before the input advances its state to the next sample.}
 \label{fig:aug_closed_loop}
\end{figure}

The condition $x^c_\tau=0$ sets $\hat s_\tau=s_\tau$, the controller
initialisation of~\cite[Remark~5.4]{kessels2025ai}, without reconstructing
$s_\tau$, because the recorded input already carries the recorded
controller's action.
Model error before $\tau$ is therefore not carried into the window.
Since $\hat y_k$ depends on the state only, the interconnection has no
algebraic loop despite $D_K\neq0$: each sample evaluates $\hat y_k$, then
$\hat u_k$ from \eqref{eq:aug_controller}, and then advances both states.

Over window starts $\mathcal T$ and window length $n_f$, the parameters
$\vartheta=(\theta_{\mathrm{base}},\theta_{\mathrm{aug}},\eta)$ minimise the
output error
\begin{equation}
 V(\vartheta)=\frac{1}{|\mathcal T|\,n_f}
 \sum_{\tau\in\mathcal T}\sum_{\ell=0}^{n_f-1}
 \norm{y_{\tau+\ell}-\hat y_{\tau+\ell|\tau}}_2^2
 \label{eq:aug_loss}
\end{equation}
subject to \eqref{eq:aug_transition} with $u_k=\hat u_k$,
\eqref{eq:aug_encoder} and \eqref{eq:aug_controller}, with the controller
fixed.
This is the truncated objective of~\cite[Eq.~(22)]{hoekstra2026lfr} with
each window simulated in closed loop
as in~\cite[Eq.~(5.12)]{kessels2025ai}; as in both, every sample of the
window is scored.

Inside the loop, the effect of a model error on the output decays with the
closed-loop modes instead of accumulating on the free axes, so $n_f$ must
span the closed-loop memory, not the plant's.
For the same reason, a small $V$ does not by itself show that the
open-loop plant is recovered.
Section~\ref{sec:setup} sets $n_f$ from this memory and gives $n$,
$n_{x_a}$, the network size and the free-run validation horizon.
